% 三2条件.tex - 工作条件

\subsection*{依托平台}

申请人将依托香港科技大学（广州）信息枢纽人工智能学域开展研究。香港科技大学（广州）是《粤港澳大湾区发展规划纲要》及《广州南沙深化面向世界的粤港澳全面合作总体方案》颁布实施以来成立的首家具有独立法人资格的内地与香港合作办学机构。学校于2022年6月正式成立，以发展融合学科为特色，探索创新人才培养模式，致力于培养面向未来的高水平创新型人才。该单位建有HPC+AI融合智算中心，理论算力达6.35P双精度（国际领先水平），支持各课题组开展学术研究；面向全校师生开放13个中央实验室，提供先进研究设备和专业技术支持，对融合学科发展发挥着至关重要的作用。该单位与70余家领军企业和知名科研机构签订了合作协议，与近10家行业龙头建立了联合实验室。丰富的产学研经验将成为本项目顺利实施的有利条件。

申请人所在学域设有专门的量子技术研究中心，总面积超过1000平方米，其中洁净室面积600余平方米，可支撑大型量子实验平台的建设与运行。申请人团队正在该中心筹建基于铷原子的大规模里德堡原子阵列实验平台，预期将为本项目的算法验证提供实验支撑。

\subsection*{计算资源}

香港科技大学（广州）建有HPC AI融合智算中心，可为本项目提供充足的计算资源：

\begin{itemize}[leftmargin=2em, topsep=0pt, itemsep=2pt]
\item \textbf{通用HPC AI平台}：6.35 PFlops（FP64）双精度浮点算力
\item \textbf{CPU资源}：Intel CPU 14656核 + AMD CPU 2560核
\item \textbf{GPU资源}：NVIDIA A800 GPU 520张 + A40 GPU 120张
\item \textbf{内存资源}：总内存容量2.3PB
\item \textbf{存储资源}：高速SSD 309TB + 大容量HDD 3.9PB
\item \textbf{国产AI平台}：19.04 PFlops（FP16）半精度算力
\end{itemize}

该计算平台可充分满足本项目大规模张量网络计算、LLM推理调用、整数规划求解等计算密集型任务的需求。特别是A800 GPU集群可支持大语言模型的本地部署与微调，确保研究数据的安全性与计算效率。

\subsection*{合作网络}

申请人与国内外研究团队和产业伙伴建立了稳定的合作关系，可为本项目提供实验验证与理论支撑：

\begin{itemize}[leftmargin=2em, topsep=0pt, itemsep=2pt]
\item \textbf{华为}：申请人与华为建立了产业合作关系，可为本项目的硬件加速器后端验证提供平台支撑。

\item \textbf{国际学术网络}：申请人曾在哈佛大学Lukin团队完成博士后训练并担任QuEra Computing顾问，在里德堡原子量子优化领域积累了深厚的国际合作基础。此外，与埼玉大学Hiroshi Shinaoka教授团队在张量网络与智能体编程方面保持合作。

\item \textbf{国内合作网络}：申请人与清华大学、中国科学院物理研究所、港科大物理系、深圳国际量子研究院等机构保持合作关系，可获得理论研究与量子模拟的支持。
\end{itemize}

\subsection*{尚缺条件与解决途径}

\begin{itemize}[leftmargin=2em, topsep=0pt, itemsep=2pt]
\item \textbf{本地里德堡原子阵列实验平台（S3）}：申请人团队正在筹建基于铷原子的实验平台，预计项目中期完成搭建。过渡期内可依托产业合作获取实验验证时间。
\item \textbf{大规模LLM调用资源（S2）}：LLM驱动的规约合成需要大量模型推理调用。\textit{解决途径}：优先使用商用API接口进行推理；若需本地化部署，可利用校内A800 GPU集群。
\item \textbf{硬件加速器平台实验时间（S3）}：量子硬件实验时间受平台可用性约束。\textit{解决途径}：通过学术合作申请实验额度，合理规划实验批次以提高利用效率；必要时以数值模拟先行验证。
\end{itemize}

